\documentclass[10pt,a4paper]{amsart}
\usepackage{amsmath,amssymb,amsthm,graphicx}
\usepackage{verbatim}
\newenvironment{reference}{\comment}{\endcomment}
%\newenvironment{reference}{}{}
\newenvironment{slogan}{\comment}{\endcomment}
\newenvironment{history}{\comment}{\endcomment}

% For commutative diagrams we use Xy-pic
\usepackage[all]{xy}

% We use 2cell for 2-commutative diagrams.
\xyoption{2cell}
% Theorem environments.
%
\theoremstyle{plain}
\newtheorem{theorem}[subsection]{Theorem}
\newtheorem{proposition}[subsection]{Proposition}
\newtheorem{lemma}[subsection]{Lemma}

\theoremstyle{definition}
\newtheorem{definition}[subsection]{Definition}
\newtheorem{example}[subsection]{Example}
\newtheorem{exercise}[subsection]{Exercise}
\newtheorem{situation}[subsection]{Situation}

\theoremstyle{remark}
\newtheorem{remark}[subsection]{Remark}
\newtheorem{remarks}[subsection]{Remarks}

\numberwithin{equation}{subsection}

% Macros
%
\def\lim{\mathop{\mathrm{lim}}\nolimits}
\def\colim{\mathop{\mathrm{colim}}\nolimits}
\def\Spec{\mathop{\mathrm{Spec}}}
\def\Hom{\mathop{\mathrm{Hom}}\nolimits}
\def\Ext{\mathop{\mathrm{Ext}}\nolimits}
\def\SheafHom{\mathop{\mathcal{H}\!\mathit{om}}\nolimits}
\def\SheafExt{\mathop{\mathcal{E}\!\mathit{xt}}\nolimits}
\def\Sch{\mathit{Sch}}
\def\Mor{\mathop{\mathrm{Mor}}\nolimits}
\def\Ob{\mathop{\mathrm{Ob}}\nolimits}
\def\Sh{\mathop{\mathit{Sh}}\nolimits}
\def\NL{\mathop{N\!L}\nolimits}
\def\CH{\mathop{\mathrm{CH}}\nolimits}
\def\proetale{{pro\text{-}\acute{e}tale}}
\def\etale{{\acute{e}tale}}
\def\QCoh{\mathit{QCoh}}
\def\Ker{\mathop{\mathrm{Ker}}}
\def\Im{\mathop{\mathrm{Im}}}
\def\Coker{\mathop{\mathrm{Coker}}}
\def\Coim{\mathop{\mathrm{Coim}}}

% Boxtimes
%
\DeclareMathSymbol{\boxtimes}{\mathbin}{AMSa}{"02}

%
% Macros for moduli stacks/spaces
%
\def\QCohstack{\mathcal{QC}\!\mathit{oh}}
\def\Cohstack{\mathcal{C}\!\mathit{oh}}
\def\Spacesstack{\mathcal{S}\!\mathit{paces}}
\def\Quotfunctor{\mathrm{Quot}}
\def\Hilbfunctor{\mathrm{Hilb}}
\def\Curvesstack{\mathcal{C}\!\mathit{urves}}
\def\Polarizedstack{\mathcal{P}\!\mathit{olarized}}
\def\Complexesstack{\mathcal{C}\!\mathit{omplexes}}
% \Pic is the operator that assigns to X its picard group, usage \Pic(X)
% \Picardstack_{X/B} denotes the Picard stack of X over B
% \Picardfunctor_{X/B} denotes the Picard functor of X over B
\def\Pic{\mathop{\mathrm{Pic}}\nolimits}
\def\Picardstack{\mathcal{P}\!\mathit{ic}}
\def\Picardfunctor{\mathrm{Pic}}
\def\Deformationcategory{\mathcal{D}\!\mathit{ef}}
\usepackage{cleveref}
\crefname{originalreference}{original reference}{original references}
\Crefname{originalreference}{Original reference}{Original references}
\makeatletter
\newlabel{lemma-category-is-colimit}{{09YU}{source}{}{sourcepointer-1}{}}
\newlabel{lemma-category-is-colimit@cref}{{[originalreference][0][]09YU}{[1][0][]source}{}{}{}}
\makeatother
\crefname{originalreference}{original reference}{original references}
\Crefname{originalreference}{Original reference}{Original references}
\makeatletter
\newlabel{lemma-constructible-maps-into-constant}{{09Z6}{source}{}{sourcepointer-2}{}}
\newlabel{lemma-constructible-maps-into-constant@cref}{{[originalreference][0][]09Z6}{[1][0][]source}{}{}{}}
\makeatother
\crefname{originalreference}{original reference}{original references}
\Crefname{originalreference}{Original reference}{Original references}
\makeatletter
\newlabel{lemma-finite-pushforward-commutes-with-base-change}{{0959}{source}{}{sourcepointer-3}{}}
\newlabel{lemma-finite-pushforward-commutes-with-base-change@cref}{{[originalreference][0][]0959}{[1][0][]source}{}{}{}}
\makeatother
\crefname{originalproposition}{original proposition}{original propositions}
\Crefname{originalproposition}{Original proposition}{Original propositions}
\makeatletter
\newlabel{limits-proposition-approximate}{{01ZA}{source}{}{sourcepointer-4}{}}
\newlabel{limits-proposition-approximate@cref}{{[originalproposition][0][]01ZA}{[1][0][]source}{}{}{}}
\makeatother
\title{Every constructible sheaf on a quasi-compact quasi-separated scheme embeds into a finite product or sum of pushforwards of constant sheaves}
\author{The Stacks Project Authors}
\date{stacks-project:etale-cohomology.tex:local-snapshot-20261011; GNU FDL 1.2 or later}
\begin{document}
\maketitle
\noindent\textit{Editorial note.} All delivered mathematical text is the author's own, copied from the pinned Stacks Project chapter file etale-cohomology.tex; no statement or proof step was rewritten, reordered or removed.\par\medskip
\noindent\textit{Editorial note.} The counted proof is the author's complete argument: the proof is written out in full: the absolute Noetherian approximation of X, the descent of the constructible sheaf to an approximant through the colimit description, the appeal to the Noetherian case, and the transport back along the pullback using base change for finite pushforwards.\par\medskip
\noindent\textit{Editorial note.} The label lemma-constructible-maps-into-constant-general names the counted statement, and the proof environment immediately following it in the source is delivered verbatim as its proof.\par\medskip
\noindent\textit{Editorial note.} The counted statement is carried by the author's own lemma environment in the chapter source.\par\medskip
\noindent\textit{Editorial note.} The statement retains the author's own bibliography reference environment to SGA4 Exposee IX Proposition 2.14, delivered verbatim as part of the counted statement.\par\medskip
\noindent\textit{Editorial note.} The preamble is the chapter's own native preamble excerpt already used for this source: the Xy-pic setup (native preamble lines 14-18) and the native theorem declarations with the mathematical macros (native preamble lines 154-217).\par\medskip
\noindent\textit{Editorial note.} Because the counted statement carries the author reference environment (SGA4 Exposee IX, Proposition 2.14), the chapter preamble excerpt includes the native verbatim package and the definitions of the reference, slogan and history environments (native preamble lines 8-13).\par\medskip
\noindent\textit{Allowed standard prerequisites (editorial).}\begin{itemize}
\item Absolute Noetherian approximation of a quasi-compact quasi-separated scheme. Limits, the proposition presenting X as a limit of schemes of finite type over Spec(Z) with affine transition morphisms. Supplies the presentation X = lim X\_t used to reduce to the Noetherian case.
\item Constructible sheaves as a colimit over the approximants. Same chapter, the lemma identifying the category of constructible sheaves on X with the colimit of the corresponding categories on the X\_t. Supplies that the given constructible sheaf is the pullback of one on some X\_t.
\item The Noetherian case and base change for finite pushforwards. Same chapter, the Noetherian form of the embedding statement and the lemma that formation of a finite pushforward commutes with base change. Supplies the embedding over X\_t and its transport to X along the pullback.
\end{itemize}\medskip
\ifdefined\hypertarget\hypertarget{sourcepointer-1}{}\fi
\noindent\textit{Source reference.} Original reference 09YU: sources/stacks-etale-cohomology/pinned-original-source.tex, source lines 12552--12552. Prerequisite pointer invoked at proof line(s) 13068; it names a standard fact of the same chapter that the counted argument uses, and every step of the counted proof built on it is retained verbatim in the delivered body.\par\medskip
\ifdefined\hypertarget\hypertarget{sourcepointer-2}{}\fi
\noindent\textit{Source reference.} Original reference 09Z6: sources/stacks-etale-cohomology/pinned-original-source.tex, source lines 12921--12921. Prerequisite pointer invoked at proof line(s) 13074; it names a standard fact of the same chapter that the counted argument uses, and every step of the counted proof built on it is retained verbatim in the delivered body.\par\medskip
\ifdefined\hypertarget\hypertarget{sourcepointer-3}{}\fi
\noindent\textit{Source reference.} Original reference 0959: sources/stacks-etale-cohomology/pinned-original-source.tex, source lines 8292--8292. Prerequisite pointer invoked at proof line(s) 13087; it names a standard fact of the same chapter that the counted argument uses, and every step of the counted proof built on it is retained verbatim in the delivered body.\par\medskip
\ifdefined\hypertarget\hypertarget{sourcepointer-4}{}\fi
\noindent\textit{Source reference.} Original proposition 01ZA: sources/stacks-chow/tags.txt, source lines 2585--2585. Prerequisite pointer invoked at proof line(s) 13065; it names a standard fact of a sibling chapter that the counted argument uses, and every step of the counted proof built on it is retained verbatim in the delivered body.\par\medskip
\begin{lemma}
\label{lemma-constructible-maps-into-constant-general}
\begin{reference}
\cite[Exposee IX, Proposition 2.14]{SGA4}
\end{reference}
Let $X$ be a quasi-compact and quasi-separated scheme.
\begin{enumerate}
\item Let $\mathcal{F}$ be a constructible sheaf of sets on $X_\etale$.
There exist an injective map of sheaves
$$
\mathcal{F} \longrightarrow
\prod\nolimits_{i = 1, \ldots, n} f_{i, *}\underline{E_i}
$$
where $f_i : Y_i \to X$ is a finite and finitely presented morphism and
$E_i$ is a finite set.
\item Let $\mathcal{F}$ be a constructible abelian sheaf on $X_\etale$.
There exist an injective map of abelian sheaves
$$
\mathcal{F} \longrightarrow
\bigoplus\nolimits_{i = 1, \ldots, n} f_{i, *}\underline{M_i}
$$
where $f_i : Y_i \to X$ is a finite and finitely presented morphism and
$M_i$ is a finite abelian group.
\item Let $\Lambda$ be a Noetherian ring.
Let $\mathcal{F}$ be a constructible sheaf of $\Lambda$-modules on $X_\etale$.
There exist an injective map of sheaves of modules
$$
\mathcal{F} \longrightarrow
\bigoplus\nolimits_{i = 1, \ldots, n} f_{i, *}\underline{M_i}
$$
where $f_i : Y_i \to X$ is a finite and finitely presented morphism and
$M_i$ is a finite $\Lambda$-module.
\end{enumerate}
\end{lemma}

\begin{proof}
We will reduce this lemma to the Noetherian case by absolute Noetherian
approximation. Namely, by
Limits, Proposition \ref{limits-proposition-approximate}
we can write $X = \lim_{t \in T} X_t$ with each $X_t$ of finite type over
$\Spec(\mathbf{Z})$ and with affine transition morphisms. By
Lemma \ref{lemma-category-is-colimit}
the category of constructible sheaves (of sets, abelian groups, or
$\Lambda$-modules) on $X_\etale$ is the colimit of the corresponding
categories for $X_t$. Thus our constructible sheaf $\mathcal{F}$
is the pullback of a similar constructible sheaf $\mathcal{F}_t$
over $X_t$ for some $t$. Then we apply the Noetherian case
(Lemma \ref{lemma-constructible-maps-into-constant})
to find an injection
$$
\mathcal{F}_t \longrightarrow
\prod\nolimits_{i = 1, \ldots, n} f_{i, *}\underline{E_i}
\quad\text{or}\quad
\mathcal{F}_t \longrightarrow
\bigoplus\nolimits_{i = 1, \ldots, n} f_{i, *}\underline{M_i}
$$
over $X_t$ for some finite morphisms $f_i : Y_i \to X_t$.
Since $X_t$ is Noetherian the morphisms $f_i$ are of finite presentation.
Since pullback is exact and since formation of $f_{i, *}$ commutes
with base change
(Lemma \ref{lemma-finite-pushforward-commutes-with-base-change}), we conclude.
\end{proof}

\end{document}
